\documentclass[11pt]{article}
\usepackage[utf8]{inputenc}
\usepackage[T1]{fontenc}
\usepackage{multicol}
\usepackage{geometry}
\usepackage{enumitem}
\usepackage{xcolor}
\usepackage{titlesec}
\usepackage{ragged2e}
\usepackage{amsmath}

\geometry{a4paper, left=1cm, right=1cm, top=0.8cm, bottom=1.2cm}

\setlength{\parindent}{0pt}
\setlength{\columnsep}{1cm}
\setlength{\columnseprule}{0.4pt}

\titleformat{\section}{\normalfont\Large\bfseries\color{blue}}{}{0em}{}

\title{\textbf{\textcolor{red}{Grandezas Diretamente e Inversamente Proporcionais - Correção}}}
\author{Professor: Jefferson}
\date{}

\begin{document}

\maketitle
\vspace{-0.5cm}

\begin{multicols}{2}

\section*{Questões e Respostas}

\begin{enumerate}[leftmargin=*, itemsep=2pt]

\item Se 5 cadernos custam R\$ 40, quanto custarão 8 cadernos? \\
\textcolor{blue}{\textbf{Resposta: R\$ 64,00}} \\
\textcolor{blue}{\textbf{Resolução:} Diretamente proporcional: mais cadernos = mais dinheiro} \\
\textcolor{blue}{$\frac{5}{8} = \frac{40}{x} \rightarrow 5x = 320 \rightarrow x = 64$}

\item Uma máquina produz 300 peças em 6 horas. Quantas peças produzirá em 10 horas? \\
\textcolor{blue}{\textbf{Resposta: 500 peças}} \\
\textcolor{blue}{\textbf{Resolução:} Diretamente proporcional: mais horas = mais peças} \\
\textcolor{blue}{$\frac{6}{10} = \frac{300}{x} \rightarrow 6x = 3000 \rightarrow x = 500$}

\item Se 12 litros de combustível permitem percorrer 180 km, quantos quilômetros serão percorridos com 20 litros? \\
\textcolor{blue}{\textbf{Resposta: 300 km}} \\
\textcolor{blue}{\textbf{Resolução:} Diretamente proporcional: mais combustível = mais km} \\
\textcolor{blue}{$\frac{12}{20} = \frac{180}{x} \rightarrow 12x = 3600 \rightarrow x = 300$}

\item Um pedreiro recebe R\$ 150 por dia. Quanto receberá em 18 dias? \\
\textcolor{blue}{\textbf{Resposta: R\$ 2.700,00}} \\
\textcolor{blue}{\textbf{Resolução:} Diretamente proporcional: mais dias = mais dinheiro} \\
\textcolor{blue}{$\frac{1}{18} = \frac{150}{x} \rightarrow x = 150 \times 18 = 2700$}

\item 4 torneiras enchem um reservatório em 2 horas. Quantas torneiras serão necessárias para encher 3 reservatórios no mesmo tempo? \\
\textcolor{blue}{\textbf{Resposta: 12 torneiras}} \\
\textcolor{blue}{\textbf{Resolução:} Diretamente proporcional: mais reservatórios = mais torneiras} \\
\textcolor{blue}{4 torneiras $\rightarrow$ 1 reservatório} \\
\textcolor{blue}{$x$ torneiras $\rightarrow$ 3 reservatórios} \\
\textcolor{blue}{$x = 4 \times 3 = 12$}

\item 7 operários constroem um muro em 5 dias. Quantos operários são necessários para construir 3 muros no mesmo período? \\
\textcolor{blue}{\textbf{Resposta: 21 operários}} \\
\textcolor{blue}{\textbf{Resolução:} Diretamente proporcional: mais muros = mais operários} \\
\textcolor{blue}{7 operários $\rightarrow$ 1 muro} \\
\textcolor{blue}{$x$ operários $\rightarrow$ 3 muros} \\
\textcolor{blue}{$x = 7 \times 3 = 21$}

\item Uma receita utiliza 250 g de farinha para fazer um bolo. Quantos gramas serão necessários para 6 bolos? \\
\textcolor{blue}{\textbf{Resposta: 1.500 g}} \\
\textcolor{blue}{\textbf{Resolução:} Diretamente proporcional: mais bolos = mais farinha} \\
\textcolor{blue}{$\frac{1}{6} = \frac{250}{x} \rightarrow x = 250 \times 6 = 1500$}

\item 3 kg de carne servem 6 pessoas. Quantos kg serão necessários para 15 pessoas? \\
\textcolor{blue}{\textbf{Resposta: 7,5 kg}} \\
\textcolor{blue}{\textbf{Resolução:} Diretamente proporcional: mais pessoas = mais carne} \\
\textcolor{blue}{$\frac{6}{15} = \frac{3}{x} \rightarrow 6x = 45 \rightarrow x = 7,5$}

\item Se 8 metros de tecido custam R\$ 200, quanto custarão 15 metros? \\
\textcolor{blue}{\textbf{Resposta: R\$ 375,00}} \\
\textcolor{blue}{\textbf{Resolução:} Diretamente proporcional: mais metros = mais dinheiro} \\
\textcolor{blue}{$\frac{8}{15} = \frac{200}{x} \rightarrow 8x = 3000 \rightarrow x = 375$}

\item Uma gráfica imprime 2.400 panfletos em 3 horas. Quantos panfletos imprimirá em 8 horas? \\
\textcolor{blue}{\textbf{Resposta: 6.400 panfletos}} \\
\textcolor{blue}{\textbf{Resolução:} Diretamente proporcional: mais horas = mais panfletos} \\
\textcolor{blue}{$\frac{3}{8} = \frac{2400}{x} \rightarrow 3x = 19200 \rightarrow x = 6400$}

\item Se 6 livros custam R\$ 90, quanto custarão 14 livros? \\
\textcolor{blue}{\textbf{Resposta: R\$ 210,00}} \\
\textcolor{blue}{\textbf{Resolução:} Diretamente proporcional: mais livros = mais dinheiro} \\
\textcolor{blue}{$\frac{6}{14} = \frac{90}{x} \rightarrow 6x = 1260 \rightarrow x = 210$}

\item Um carro percorre 150 km com 10 litros de combustível. Quantos quilômetros percorrerá com 25 litros? \\
\textcolor{blue}{\textbf{Resposta: 375 km}} \\
\textcolor{blue}{\textbf{Resolução:} Diretamente proporcional: mais combustível = mais km} \\
\textcolor{blue}{$\frac{10}{25} = \frac{150}{x} \rightarrow 10x = 3750 \rightarrow x = 375$}

\item 9 trabalhadores constroem uma cerca em 4 dias. Quantos trabalhadores serão necessários para construir 3 cercas no mesmo período? \\
\textcolor{blue}{\textbf{Resposta: 27 trabalhadores}} \\
\textcolor{blue}{\textbf{Resolução:} Diretamente proporcional: mais cercas = mais trabalhadores} \\
\textcolor{blue}{9 trabalhadores $\rightarrow$ 1 cerca} \\
\textcolor{blue}{$x$ trabalhadores $\rightarrow$ 3 cercas} \\
\textcolor{blue}{$x = 9 \times 3 = 27$}

\item Uma fábrica produz 500 garrafas em 5 horas. Quantas garrafas produzirá em 12 horas? \\
\textcolor{blue}{\textbf{Resposta: 1.200 garrafas}} \\
\textcolor{blue}{\textbf{Resolução:} Diretamente proporcional: mais horas = mais garrafas} \\
\textcolor{blue}{$\frac{5}{12} = \frac{500}{x} \rightarrow 5x = 6000 \rightarrow x = 1200$}

\item Se 4 kg de arroz custam R\$ 32, quanto custarão 10 kg? \\
\textcolor{blue}{\textbf{Resposta: R\$ 80,00}} \\
\textcolor{blue}{\textbf{Resolução:} Diretamente proporcional: mais kg = mais dinheiro} \\
\textcolor{blue}{$\frac{4}{10} = \frac{32}{x} \rightarrow 4x = 320 \rightarrow x = 80$}

\item Uma impressora imprime 120 páginas em 6 minutos. Quantas páginas imprimirá em 15 minutos? \\
\textcolor{blue}{\textbf{Resposta: 300 páginas}} \\
\textcolor{blue}{\textbf{Resolução:} Diretamente proporcional: mais minutos = mais páginas} \\
\textcolor{blue}{$\frac{6}{15} = \frac{120}{x} \rightarrow 6x = 1800 \rightarrow x = 300$}

\item 2 caixas contêm 48 lápis. Quantos lápis haverá em 7 caixas iguais? \\
\textcolor{blue}{\textbf{Resposta: 168 lápis}} \\
\textcolor{blue}{\textbf{Resolução:} Diretamente proporcional: mais caixas = mais lápis} \\
\textcolor{blue}{$\frac{2}{7} = \frac{48}{x} \rightarrow 2x = 336 \rightarrow x = 168$}

\item Um caminhão transporta 8 toneladas de areia em uma viagem. Quantas toneladas transportará em 15 viagens? \\
\textcolor{blue}{\textbf{Resposta: 120 toneladas}} \\
\textcolor{blue}{\textbf{Resolução:} Diretamente proporcional: mais viagens = mais toneladas} \\
\textcolor{blue}{$\frac{1}{15} = \frac{8}{x} \rightarrow x = 8 \times 15 = 120$}

\item 5 metros de fio pesam 200 g. Quanto pesarão 18 metros do mesmo fio? \\
\textcolor{blue}{\textbf{Resposta: 720 g}} \\
\textcolor{blue}{\textbf{Resolução:} Diretamente proporcional: mais metros = mais peso} \\
\textcolor{blue}{$\frac{5}{18} = \frac{200}{x} \rightarrow 5x = 3600 \rightarrow x = 720$}

\item Se 3 litros de tinta pintam 24 m² de parede, quantos metros quadrados serão pintados com 10 litros? \\
\textcolor{blue}{\textbf{Resposta: 80 m²}} \\
\textcolor{blue}{\textbf{Resolução:} Diretamente proporcional: mais tinta = mais área} \\
\textcolor{blue}{$\frac{3}{10} = \frac{24}{x} \rightarrow 3x = 240 \rightarrow x = 80$}

\item Três trabalhadores fazem um serviço em 12 dias. Quantos dias levariam 6 trabalhadores para fazer o mesmo serviço? \\
\textcolor{blue}{\textbf{Resposta: 6 dias}} \\
\textcolor{blue}{\textbf{Resolução:} Inversamente proporcional: mais trabalhadores = menos dias} \\
\textcolor{blue}{$3 \times 12 = 6 \times x \rightarrow 36 = 6x \rightarrow x = 6$}

\item Se 4 máquinas fazem um trabalho em 5 horas, quantas horas seriam necessárias para que 10 máquinas realizassem o mesmo trabalho? \\
\textcolor{blue}{\textbf{Resposta: 2 horas}} \\
\textcolor{blue}{\textbf{Resolução:} Inversamente proporcional: mais máquinas = menos horas} \\
\textcolor{blue}{$4 \times 5 = 10 \times x \rightarrow 20 = 10x \rightarrow x = 2$}

\item Dez operários fazem um trabalho em 8 dias. Quantos operários são necessários para que termine em 4 dias? \\
\textcolor{blue}{\textbf{Resposta: 20 operários}} \\
\textcolor{blue}{\textbf{Resolução:} Inversamente proporcional: mais operários = menos dias} \\
\textcolor{blue}{$10 \times 8 = x \times 4 \rightarrow 80 = 4x \rightarrow x = 20$}

\item Uma torneira enche um reservatório em 10 horas. Se abrirmos mais 2 torneiras iguais, quanto tempo será preciso? \\
\textcolor{blue}{\textbf{Resposta: 3,33 horas ou 3h20min}} \\
\textcolor{blue}{\textbf{Resolução:} Inversamente proporcional: mais torneiras = menos tempo} \\
\textcolor{blue}{$1 \times 10 = 3 \times x \rightarrow 10 = 3x \rightarrow x \approx 3,33$}

\item Um grupo de 8 pessoas leva 4 dias para pintar uma casa. Quantos dias levará um grupo de 12 pessoas? \\
\textcolor{blue}{\textbf{Resposta: 2,67 dias ou 2 dias e 16h}} \\
\textcolor{blue}{\textbf{Resolução:} Inversamente proporcional: mais pessoas = menos dias} \\
\textcolor{blue}{$8 \times 4 = 12 \times x \rightarrow 32 = 12x \rightarrow x \approx 2,67$}

\item Uma equipe de 5 jardineiros leva 15 horas para terminar um jardim. Quantas horas levariam 10 jardineiros? \\
\textcolor{blue}{\textbf{Resposta: 7,5 horas}} \\
\textcolor{blue}{\textbf{Resolução:} Inversamente proporcional: mais jardineiros = menos horas} \\
\textcolor{blue}{$5 \times 15 = 10 \times x \rightarrow 75 = 10x \rightarrow x = 7,5$}

\item Um serviço é feito por 9 trabalhadores em 20 dias. Quantos trabalhadores são necessários para que o serviço seja feito em 15 dias? \\
\textcolor{blue}{\textbf{Resposta: 12 trabalhadores}} \\
\textcolor{blue}{\textbf{Resolução:} Inversamente proporcional: menos dias = mais trabalhadores} \\
\textcolor{blue}{$9 \times 20 = x \times 15 \rightarrow 180 = 15x \rightarrow x = 12$}

\item 4 carros levam 18 horas para completar um trabalho. Quantos carros seriam necessários para completar em 6 horas? \\
\textcolor{blue}{\textbf{Resposta: 12 carros}} \\
\textcolor{blue}{\textbf{Resolução:} Inversamente proporcional: menos horas = mais carros} \\
\textcolor{blue}{$4 \times 18 = x \times 6 \rightarrow 72 = 6x \rightarrow x = 12$}

\item Um encanador leva 9 horas para consertar 3 canos. Quantas horas levaria para consertar 6 canos? \\
\textcolor{blue}{\textbf{Resposta: 18 horas}} \\
\textcolor{blue}{\textbf{Resolução:} Diretamente proporcional: mais canos = mais horas} \\
\textcolor{blue}{$\frac{3}{6} = \frac{9}{x} \rightarrow 3x = 54 \rightarrow x = 18$}

\item Se 15 pintores fazem um trabalho em 5 dias, quantos dias levariam 30 pintores? \\
\textcolor{blue}{\textbf{Resposta: 2,5 dias}} \\
\textcolor{blue}{\textbf{Resolução:} Inversamente proporcional: mais pintores = menos dias} \\
\textcolor{blue}{$15 \times 5 = 30 \times x \rightarrow 75 = 30x \rightarrow x = 2,5$}

\end{enumerate}

\end{multicols}

\end{document}
